\documentclass[11pt]{article}
\setlength{\topmargin}{-0.3in}
\setlength{\oddsidemargin}{-0.1in}
\setlength{\evensidemargin}{-0.1in}
\setlength{\textwidth}{6.7in}
\setlength{\headheight}{0in}
\setlength{\headsep}{0in}
\setlength{\topskip}{0.5in}
\setlength{\textheight}{9.25in}
\setlength{\parindent}{0.0in}
\setlength{\parskip}{1em}
\usepackage{common,graphicx,hyperref,epsfig}

\pagestyle{empty}

\begin{document}

\today

Prof. Reggie Madjoe\\
Chair, Department of Physics\\
University of the Western Cape\\
Private Bag X17\\
7535 Bellville\\
Western Cape Province\\
South Africa

Dear Prof. Madjoe:

I am writing to apply for the position of lecturer in the Astrophysics
group within the Physics Deparment. My area of research expertise is 

I earned my Ph.D. from Michigan
State University in 2008

am a graduate
student at Prestigious University working on a dissertation under the
direction of Professor Prominent Figure. Currently revising the third
of five chapters, I expect to complete all work for the Ph.D. by May
of 1999. I believe that my teaching and tutoring experience combined
with my course work and research background in rhetoric and
composition theory make me a strong candidate for the position
outlined in your notice.

As my curriculum vitae shows, I have had excellent opportunities to
teach a variety of writing courses during my graduate studies,
including developmental writing, first-year writing for both native
speakers and second language students, advanced writing, and business
writing. I have also worked as a teaching mentor for new graduate
students, a position that involved instruction in methods of
composition teaching, development of course materials, and evaluation
of new graduate instructors. Among the most satisfying experiences for
me as a teacher has been instructing students on an individual basis
as a tutor in our university Writing Lab. Even as a classroom
instructor, I find that I always look forward to the individual
conferences that I hold with my students several times during the
semester because I believe this kind of one-on-one interaction to be
essential to their development as writers.

My work in the composition classroom has provided me with the
inspiration as well as a kind of laboratory for my dissertation
research. My project, The I Has It: Applications of Recent Models of
Subjectivity in Composition Theory, examines the shift since the 1960s
from expressive models of writing toward now-dominant postmodern
conceptions of decentered subjectivity and self-construction through
writing.  I argue that these more recent theoretical models, while
promising, cannot have the liberating effects that are claimed for
them without a concomitant reconception of writing pedagogy and the
dynamics of the writing classroom. I relate critical readings of
theoretical texts to my own pedagogical experiments as a writing
teacher, using narratives of classroom successes and failures as the
bases for critical reflection on postmodern composition theory. After
developing my dissertation into a book manuscript, I plan to continue
my work in current composition theory through a critical examination
of the rhetoric of technological advancement in the computer-mediated
writing classroom.

My interest in the computer classroom has grown out of recent
experience teaching composition in that environment.  In these courses
my students have used computers for writing and turning in notes and
essays, communicating with one another and with me, conducting library
catalogue research and web research, and creating websites. I have
encouraged my students to think and write critically about their
experiences with technology, both in my class and elsewhere, even as
we have used technology to facilitate our work in the course. Syllabi
and other materials for my writing courses can be viewed at my
website: <http://machine.prestigious.edu/~name>. In all of my writing
courses I encourage students to become critical readers, thinkers, and
writers; my goal is always not only to promote their intellectual
engagement with cultural texts of all kinds but also to help them
become more discerning readers of and forceful writers about the world
around them.

I have included my curriculum vitae and would be happy to send you
additional materials such as a dossier of letters of reference,
writing samples, teaching evaluations, and past and proposed course
syllabi. I will be available to meet with you for an interview at
either the MLA or the CCCC convention, or elsewhere at your
convenience. I can be reached at my home phone number before December
19; between then and the start of the MLA convention, you can reach me
at (123) 456-7890. I thank you for your consideration and look forward
to hearing from you.

Sincerely,\\
\begin{minipage}{7.5in}
\epsfig{file=signature.ps,height=0.75in}
\end{minipage}
Kenneth Cavagnolo

\end{document}
